\Article{9}

Secundo loco considerabimus progressionem hancce

\[
\begin{aligned}
&1+x^{\frac12}\frac{1-x^m}{1-x}+x\frac{(1-x^m)(1-x^{m-1})}{(1-x)(1-xx)}\\
&\qquad+x^{\frac32}\frac{(1-x^m)(1-x^{m-1})(1-x^{m-2})}{(1-x)(1-xx)(1-x^3)}+\text{etc.}
\end{aligned}
\]

sive

\[
\begin{aligned}
1+x^{\frac12}(m,1)+x(m,2)+x^{\frac32}(m,3)+xx(m,4)+\text{etc.}
\end{aligned}
\]

quam per $F(x,m)$ denotabimus.\quad Restringemus hanc disquisitionem ad casum eum, ubi $m$ est integer positivus, ita ut haec quoque series semper abrumpatur 
\SourceLeaf{31}

cum termino $m+1^{\mathrm{to}}$, qui est $=x^{\frac12m}(m,m)$.\quad Quum sit.

\[
\begin{aligned}
(m,m)=1,\quad(m,m-1)=(m,1),\quad(m,m-2)=(m,2),\;\text{etc.}
\end{aligned}
\]

progressio ita quoque exhiberi poterit:

\[
\begin{aligned}
F(x,m)&=x^{\frac12m}+x^{\frac12(m-1)}(m,1)\\&\qquad+x^{\frac12(m-2)}(m,2)+x^{\frac12(m-3)}(m,3)+\text{etc.}
\end{aligned}
\]

Hinc fit

\[
\begin{aligned}
(1+x^{\frac12m+\frac12})F(x,m)&=1+x^{\frac12}(m,1)+x(m,2)\\&\qquad+x^{\frac32}(m,3)+\text{etc.}\\
&+x^{\frac12}\pdot x^m+x\pdot x^{m-1}(m,1)\\&\qquad+x^{\frac32}\pdot x^{m-2}(m,2)+\text{etc.}
\end{aligned}
\]

Quare quum habeatur (art. 5. II)

\[
\begin{aligned}
(m,1)+x^m&=(m+1,1)\\
(m,2)+x^{m-1}(m,1)&=(m+1,2)\\
(m,3)+x^{m-2}(m,2)&=(m+1,3),\;\text{etc.}
\end{aligned}
\]

provenit

\[
\begin{aligned}
(1+x^{\frac12m+\frac12})F(x,m)=F(x,m+1)\ .\; .\; .\; .\; .\; .\; .\ [3]
\end{aligned}
\]

Sed fit $F(x,0)=1$: quamobrem erit

\[
\begin{aligned}
F(x,1)&=1+x^{\frac12}\\
F(x,2)&=(1+x^{\frac12})(1+x)\\
F(x,3)&=(1+x^{\frac12})(1+x)(1+x^{\frac32}),\;\text{etc.}
\end{aligned}
\]

sive generaliter

\[
\begin{aligned}
F(x,m)&=(1+x^{\frac12})(1+x)(1+x^{\frac32})\\&\qquad\dotsfour(1+x^{\frac12m})\ .\; .\; .\; .\quad[4]
\end{aligned}
\]

\Article{10}

Praemissis hisce disquisitionibus praeliminaribus iam propius ad propositum nostrum accedamus.\quad Quum pro valore primo ipsius $n$ quadrata $1$, $4$, $9\dotsfour(\frac12(n-1))^2$ omnia inter se incongrua sint secundum modulum $n$, patet, illorum residua minima secundum hunc modulum cum numeris $a$ identica esse debere, adeoque

\[
\begin{aligned}
\Sigma\cos ak\omega&=\cos k\omega+\cos4k\omega+\cos9k\omega\\&\qquad+\text{etc.}+\cos(\tfrac12(n-1))^2k\omega\\
\Sigma\sin ak\omega&=\sin k\omega+\sin4k\omega+\sin9k\omega\\&\qquad+\text{etc.}+\sin(\tfrac12(n-1))^2k\omega
\end{aligned}
\]

Perinde quum eadem quadrata $1$, $4$, $9\dotsfour(\frac12(n-1))^2$ ordine inverso congrua sint his $(\frac12(n+1))^2$, $(\frac12(n+3))^2$, $(\frac12(n+5))^2\dotsfive(n-1)^2$, etiam erit

\SourceLeaf{32}

\[
\begin{aligned}
\Sigma\cos ak\omega&=\cos(\tfrac12(n+1))^2k\omega\\&\qquad+\cos(\tfrac12(n+3))^2k\omega+\text{etc.}+\cos(n-1)^2k\omega\\
\Sigma\sin ak\omega&=\sin(\tfrac12(n+1))^2k\omega\\&\qquad+\sin(\tfrac12(n+3))^2k\omega+\text{etc.}+\sin(n-1)^2k\omega
\end{aligned}
\]

Statuendo itaque

\[
\begin{aligned}
T&=1+\cos k\omega+\cos4k\omega+\cos9k\omega\\&\qquad+\text{etc.}+\cos(n-1)^2k\omega\\
U&=\phantom{1+}\sin k\omega+\sin4k\omega+\sin9k\omega\\&\qquad+\text{etc.}+\sin(n-1)^2k\omega
\end{aligned}
\]

erit

\[
\begin{aligned}
1+2\Sigma\cos ak\omega&=T\\
2\Sigma\sin ak\omega&=U
\end{aligned}
\]

Hinc patet, summationes, quales in art. 1. propositae sunt, pendere a summatione serierum $T$ et $U$, quocirca, missis illis, disquisitionem nostram his adaptabimus, eaque generalitate absolvemus, ut non modo valores primos ipsius $n$, sed quoscunque compositos complectatur.\quad Numerum $k$ autem supponemus ad $n$ primum esse: nullo enim negotio casus is, ubi $k$ et $n$ divisorem communem haberent, ad hunc reduci poterit.

\Article{11}

Designemus quantitatem imaginariam $\surd{-1}$ per $i$, statuamusque

\[
\begin{aligned}
\cos k\omega+i\sin k\omega=r
\end{aligned}
\]

unde erit $r^n=1$, sive $r$ radix aequationis $r^n-1=0$.\quad Facile perspicietur, omnes numeros $k$, $2k$, $3k\dotsfour(n-1)k$ per $n$ non divisibiles atque inter se secundum modulum $n$ incongruos esse: hinc potestates ipsius $r$

\[
\begin{aligned}
1,\;r,\;rr,\;r^3\dotsfour r^{n-1}
\end{aligned}
\]

omnes erunt inaequales, singulae vero quoque aequationi $x^n-1=0$ satisfacient. Hanc ob caussam hae potestates \emph{omnes} radices aequationis $x^n-1=0$ repraesentabunt.

Hae conclusiones non valerent, si $k$ divisorem communem haberet cum $n$. Si enim $\nu$ esset talis divisor communis, foret $k\pdot\dfrac n\nu$ per $n$ divisibilis, adeoque potestas inferior quam $r^n$, puta $r^{\frac n\nu}$, unitati aequalis.\quad In hoc itaque casu potestates ipsius $r$ ad summum $\dfrac n\nu$ radices aequationis $x^n-1=0$ exhibebunt, et quidem revera tot radices diversas sistent, si $\nu$ est divisor communis \emph{maximus} nume\SourceLeaf{33}

rorum $k$, $n$.\quad In casu nostro, ubi $k$ et $n$ supponuntur inter se primi, $r$ commode dici potest \emph{radix propria} aequationis $x^n-1=0$: contra in casu altero, ubi $k$ et $n$ haberent divisorem communem (maximum) $\nu$, $r$ vocaretur \emph{radix impropria} illius aequationis, manifesto autem tunc eadem $r$ foret radix propria aequationis $x^{\frac n\nu}-1=0$.\quad Radix impropria simplicissima est unitas, in eoque casu, ubi $n$ est numerus primus, impropriae aliae omnino non dabuntur.

\Article{12}

Quodsi iam statuimus

\[
\begin{aligned}
W=1+r+r^4+r^9+\text{etc.}+\text{etc.}+r^{(n-1)^2}
\end{aligned}
\]

patet fieri $W=T+iU$, adeoque $T$ esse partem realem ipsius $W$, atque $U$ prodire ex parte imaginaria ipsius $W$ factore $i$ suppresso.\quad Totum itaque negotium reducitur ad inventionem summae $W$: ad hunc finem vel series in art. 6 considerata, vel ea quam in art. 9 summare docuimus, adhiberi potest, prior tamen minus idonea est in casu eo, ubi $n$ est numerus par.\quad Nihilominus lectoribus gratum fore speramus, si casum eum, ubi $n$ impar est, secundum methodum duplicem tractemus.

Supponamus itaque primo, $n$ esse numerum imparem, $r$ designare radicem propriam aequationis $x^n-1=0$ quamcunque, et in functione $f(x,m)$ statui $x=r$, atque $m=n-1$.\quad Hinc patet fieri

\[
\begin{aligned}
\frac{1-x^m}{1-x}&=\frac{1-r^{-1}}{1-r}=-r^{-1}\\
\frac{1-x^{m-1}}{1-xx}&=\frac{1-r^{-2}}{1-rr}=-r^{-2}\\
\frac{1-x^{m-2}}{1-x^3}&=\frac{1-r^{-3}}{1-r^3}=-r^{-3},\;\text{etc.}
\end{aligned}
\]

usque ad

\[
\begin{aligned}
\frac{1-x}{1-x^m}=\frac{1-r^{-m}}{1-r^m}=-r^{-m}
\end{aligned}
\]

(Haud superfluum erit monere, has aequationes eatenus tantum valere, quatenus $r$ supponitur radix propria: si enim esset $r$ radix impropria, in quibusdam illarum fractionum numerator et denominator simul evanescerent, adeoque fractiones indeterminatae fierent).

Hinc deducimus aequationem sequentem

\SourceLeaf{34}

\[
\begin{aligned}
f(r,n-1)&=1+r^{-1}+r^{-3}+r^{-6}+\text{etc.}+r^{-\frac12(n-1)n}\\
&=(1-r)(1-r^3)(1-r^5)\dotsfive(1-r^{n-2})
\end{aligned}
\]

Eadem aequatio etiamnum valebit, si pro $r$ substituitur $r^\lambda$, designante $\lambda$ integrum quemcunque ad $n$ primum: tunc enim etiam $r^\lambda$ erit radix propria aequationis $x^n-1=0$.\quad Scribamus itaque pro $r$, $r^{n-2}$ sive quod idem est $r^{-2}$, eritque

\[
\begin{aligned}
&1+r^2+r^6+r^{12}+\text{etc.}+r^{(n-1)n}\\&\qquad=(1-r^{-2})(1-r^{-6})(1-r^{-10})\ldots(1-r^{-2(n-2)})
\end{aligned}
\]

Multiplicemus utramque partem huius aequationis per

\[
\begin{aligned}
r\pdot r^3\pdot r^5\dotsfour r^{(n-2)}=r^{\frac14(n-1)^2}
\end{aligned}
\]

prodibitque, propter

\[
\begin{aligned}
r^{2+\frac14(n-1)^2}&=r^{\frac14(n-3)^2},\\r^{(n-1)n+\frac14\,n-1)^2}=r^{\frac14(n+1)^2}\\
r^{6+\frac14(n-1)^2}&=r^{\frac14(n-5)^2},\\r^{(n-2)(n-1)+\frac14(n-1)^2}=r^{\frac14(n+3)^2}\\
r^{12+\frac14(n-1)^2}&=r^{\frac14(n-7)^2},\\r^{(n-3)(n-2)+\frac14(n-1)^2}=r^{\frac14(n+5)^2},\;\text{etc.}
\end{aligned}
\]

aequatio sequens

\[
\begin{aligned}
&r^{\frac14(n-1)^2}+r^{\frac14(n-3)^2}+r^{\frac14(n-5)^2}+\text{etc.}+r+1\\
&+r^{\frac14(n+1)^2}+r^{\frac14(n+3)^2}+r^{\frac14(n+5)^2}\\&\qquad+\text{etc.}+r^{\frac14(2n-2)^2}\\
&\qquad=(r-r^{-1})(r^3-r^{-3})(r^5-r^{-5})\dotsfour(r^{n-2}-r^{-n+2})
\end{aligned}
\]

aut, partibus membri primi aliter dispositis,

\[
\begin{aligned}
&1+r+r^4+\text{etc.}+r^{(n-1)^2}\\&\qquad=(r-r^{-1})(r^3-r^{-3})\ldots(r^{n-2}-r^{-n+2})\quad\dotstwo\quad[5]
\end{aligned}
\]
